La factibilidad de recursos comprende el análisis completo de todos los recursos necesarios (humanos, materiales, financieros, tecnológicos y de tiempo) para la puesta en marcha de un proyecto  con éxito. Podría decirse que este apartado funciona más como un resumen de todos los análisis elaborados previamente. 

En esencia el análisis de factibilidad de recursos abarca todos los anteriores análisis elaborados durante la sección \ref{sec:analisis_factibilidad}. Este estudio comprende la factibilidad técnica, económica, operativa, legal y de tiempo, campos los cuales ya fueron desarrollados con anterioridad, demostrando que el proyecto es factible y alcanzable considerando todos los tipos de recursos disponibles analizados. 

A partir de los análisis previos, se determina que el proyecto es \textbf{altamente factible} en materia de recursos, fundamentándose en los siguientes tres pilares:

\begin{itemize}
    \item \textbf{Recursos humanos y temporales:} El desarrollo del prototipo recae exclusivamente en los tres autores del presente Trabajo Terminal, quienes fungen como desarrolladores, analistas y gestores del sistema bajo la metodología Scrumban. Como se demostró en la factibilidad operativa, el equipo posee el perfil técnico necesario (dominio de programación móvil, bases de datos espaciales y diseño algorítmico) y cuenta con la disponibilidad de tiempo requerida dentro de los plazos establecidos por el calendario académico de la ESCOM.
    
    \item \textbf{Recursos tecnológicos:} El proyecto adopta una arquitectura de software basada primordialmente en tecnologías de código abierto (Flutter, Dart, PostgreSQL, PostGIS, pgRouting y datos cartográficos de OpenStreetMap bajo licencia ODbL) y plataformas bajo el modelo de \textit{Backend as a Service} (Supabase) en sus capas gratuitas. Asimismo, las cuotas de procesamiento para las APIs comerciales (Google Maps Platform y Qwen3.7-Flash) representan un volumen manejable que el equipo puede absorber directamente, mientras que los servicios externos de operación (mapa base, búsqueda de lugares, correo y notificaciones \textit{push}) se cubren con planes gratuitos, con la única condición de registrar un dominio propio para el envío de correos. No se requiere la adquisición de licencias de software privativo ni la renta de servidores dedicados de alto costo.
    
    \item \textbf{Recursos materiales y de infraestructura:} Las actividades de diseño, codificación y preprocesamiento de datos se llevarán a cabo utilizando los equipos de cómputo personales de los integrantes del equipo. De igual forma, los servicios básicos requeridos para la operatividad del desarrollo, tales como el suministro de energía eléctrica y el acceso a conexiones de internet de banda ancha, ya se encuentran cubiertos en los espacios de trabajo domiciliarios de los desarrolladores, eliminando la necesidad de arrendar oficinas o adquirir mobiliario especial.

    \item\textbf{Recursos financieros y técnicos}: El cálculo de sueldos y el valor de mercado de este trabajo especializado (roles de trabajo) asciende a \$232,797.60 MXN, al ser elaborado por estudiantes como requisito de titulación, no representa un desembolso económico real. El uso de servicios y software de terceros se contempla mayormente mediante licencias gratuitas y adicionalmente, el equipo cuenta con hardware propio que cumple con los requisitos técnicos identificados (mínimo de 8 GB a 16 GB de RAM para la virtualización mediante Android Studio), eliminando la necesidad de adquirir nuevo equipo de cómputo.

    La evaluación económica concluye que, si bien el proyecto tiene una valuación total estimada en \$273,406.78 MXN (incluyendo imprevistos y el costo de los recursos materiales preexistentes como internet y energía eléctrica), el egreso de capital efectivo es sumamente reducido. El desembolso real proyectado se limita a \$3,486.46 MXN, equivalente a un esfuerzo mensual aproximado de \$116.22 MXN por integrante durante los diez meses de desarrollo. Este gasto mínimo cubrirá de manera holgada las cuotas excedentes de las APIs de preprocesamiento (a lo mucho 10,500 peticiones a Google Street View Static API y tokens del modelo Qwen3.7-Flash), así como las suscripciones de apoyo operativo (Overleaf y Udemy).
\end{itemize}

En conclusión, la disponibilidad garantizada de las herramientas de software, la suficiencia del hardware propio y la capacidad técnica del equipo de trabajo aseguran que no existen impedimentos o cuellos de botella a nivel de recursos que pongan en riesgo la materialización del prototipo dentro del ciclo escolar programado.